\documentclass{article}

\usepackage{amsmath}
\usepackage{amssymb}
\usepackage{amsthm}
\usepackage{enumitem}
\usepackage[polish]{babel}
\usepackage[utf8]{inputenc}
\usepackage[T1]{fontenc}

\title{Przestrzenie metryczne, unormowane i Hilberta}
\date{\today}

\newtheorem{definicja}{Definicja}
\newtheorem{twierdzenie}{Twierdzenie}
\newtheorem{przyklad}{Przykład}

\begin{document}
\maketitle

\section{Przestrzeń metryczna}

\begin{definicja}
Niech $X$ będzie niepustym zbiorem. Metryką na zbiorze $X$ nazywamy funkcję
\[
d:X\times X\rightarrow[0,\infty),
\]
która dla wszystkich $x,y,z\in X$ spełnia warunki:
\begin{enumerate}[label=(M\arabic*)]
    \item $d(x,y)=0 \iff x=y$,
    \item $d(x,y)=d(y,x)$,
    \item $d(x,z)\leq d(x,y)+d(y,z)$ -- nierówność trójkąta.
\end{enumerate}
Parę $(X,d)$ nazywamy przestrzenią metryczną.
\end{definicja}

Z warunków metryki wynika, że odległość dowolnych dwóch punktów jest zawsze liczbą nieujemną. Przestrzeń metryczną tworzy więc zbiór wraz z~pojęciem odległości.

\subsection{Kula i kule otwarte}
Dla $x_0\in X$ i $r>0$ definiujemy:
\[
K(x_0,r)=\{x\in X:d(x,x_0)<r\}
\]
-- kulę otwartą o środku $x_0$ i promieniu $r$, oraz
\[
\overline{K}(x_0,r)=\{x\in X:d(x,x_0)\leq r\}
\]
-- kulę domkniętą.

Kule są podstawą definicji zbieżności i ciągłości w przestrzeniach metrycznych.

\subsection{Przykłady metryk}
\begin{przyklad}
Na $\mathbb R^n$ można określić metrykę euklidesową:
\[
d_2(x,y)=\sqrt{\sum_{i=1}^{n}(x_i-y_i)^2}.
\]
\end{przyklad}

\begin{przyklad}
Dla $x,y\in\mathbb R^n$ ważne są także metryki:
\[
d_1(x,y)=\sum_{i=1}^{n}|x_i-y_i|,
\qquad
d_\infty(x,y)=\max_{1\leq i\leq n}|x_i-y_i|.
\]
\end{przyklad}

\begin{przyklad}
Na dowolnym niepustym zbiorze $X$ metrykę można określić wzorem
\[
d(x,y)=
\begin{cases}
0,&x=y,\\
1,&x\neq y.
\end{cases}
\]
Nazywa się ją metryką dyskretną.
\end{przyklad}

\begin{przyklad}
W przestrzeni funkcji ciągłych $C[a,b]$ jedną z metryk jest
\[
d_\infty(f,g)=\max_{t\in[a,b]}|f(t)-g(t)|.
\]
\end{przyklad}

\subsection{Zbieżność i ciąg Cauchy'ego}
Mówimy, że ciąg $(x_n)\subset X$ jest zbieżny do $x\in X$, gdy
\[
\forall_{\varepsilon>0}\exists_{N}\forall_{n\geq N}\ d(x_n,x)<\varepsilon.
\]

Ciąg $(x_n)$ nazywamy ciągiem Cauchy'ego, jeżeli
\[
\forall_{\varepsilon>0}\exists_{N}\forall_{m,n\geq N}\ d(x_m,x_n)<\varepsilon.
\]

Każdy ciąg zbieżny jest ciągiem Cauchy'ego, natomiast odwrotna implikacja nie musi zachodzić.

\begin{definicja}
Przestrzeń metryczną $(X,d)$ nazywamy zupełną, jeżeli każdy ciąg Cauchy'ego elementów tej przestrzeni jest zbieżny do pewnego elementu należącego do $X$.
\end{definicja}

\begin{przyklad}
Przestrzenie $\mathbb R^n$ i $\mathbb C^n$ z metryką euklidesową są zupełne.
Natomiast przedział $(0,1)$ z metryką zwykłą nie jest zupełny, ponieważ ciąg $x_n=1/n$ jest ciągiem Cauchy'ego, ale nie ma granicy należącej do $(0,1)$.
\end{przyklad}

\section{Przestrzeń unormowana}

\begin{definicja}
Niech $V$ będzie przestrzenią liniową nad ciałem $\mathbb R$ albo $\mathbb C$. Normą nazywamy funkcję
\[
\|\cdot\|:V\rightarrow[0,\infty),
\]
która dla wszystkich $x,y\in V$ i wszystkich skalarów $\alpha$ spełnia:
\begin{enumerate}[label=(N\arabic*)]
    \item $\|x\|=0 \iff x=0$,
    \item $\|\alpha x\|=|\alpha|\|x\|$,
    \item $\|x+y\|\leq\|x\|+\|y\|$ -- nierówność trójkąta.
\end{enumerate}
Parę $(V,\|\cdot\|)$ nazywamy przestrzenią unormowaną.
\end{definicja}

Każda norma wyznacza metrykę wzorem
\[
d(x,y)=\|x-y\|.
\]
Odległość taka nazywa się metryką indukowaną przez normę. Nie każda metryka pochodzi od normy -- na przykład metryka dyskretna na przestrzeni liniowej nie jest indukowana przez normę.

\subsection{Przykłady norm}

\begin{przyklad}
Na $\mathbb R^n$ lub $\mathbb C^n$ mamy normy
\[
\|x\|_1=\sum_{i=1}^n|x_i|,
\qquad
\|x\|_2=\sqrt{\sum_{i=1}^n|x_i|^2},
\qquad
\|x\|_\infty=\max_{1\leq i\leq n}|x_i|.
\]
\end{przyklad}

\begin{przyklad}
W przestrzeni $C[a,b]$ wzór
\[
\|f\|_\infty=\max_{t\in[a,b]}|f(t)|
\]
określa normę. Odpowiadająca jej metryka to
\[
d_\infty(f,g)=\|f-g\|_\infty=\max_{t\in[a,b]}|f(t)-g(t)|.
\]
\end{przyklad}

\begin{przyklad}
Dla ciągu $x=(x_n)$ przestrzeń $\ell^2$ składa się z ciągów spełniających
\[
\sum_{n=1}^{\infty}|x_n|^2<\infty.
\]
Normą jest
\[
\|x\|_2=\left(\sum_{n=1}^{\infty}|x_n|^2\right)^{1/2}.
\]
\end{przyklad}

\subsection{Przestrzeń Banacha}
\begin{definicja}
Przestrzeń unormowaną zupełną względem metryki indukowanej przez normę nazywamy przestrzenią Banacha.
\end{definicja}

Przykłady przestrzeni Banacha:
\begin{itemize}
    \item $\mathbb R^n$ i $\mathbb C^n$ z dowolną z norm $\|\cdot\|_1$, $\|\cdot\|_2$ lub $\|\cdot\|_\infty$,
    \item $C[a,b]$ z normą $\|\cdot\|_\infty$,
    \item $\ell^2$ z normą $\|\cdot\|_2$,
    \item przestrzenie $L^p$, $1\leq p<\infty$, z normą
    \[
    \|f\|_p=\left(\int |f|^p\,d\mu\right)^{1/p}.
    \]
\end{itemize}

Przestrzeń $C[a,b]$ z normą
\[
\|f\|_1=\int_a^b|f(t)|\,dt
\]
nie jest zupełna.

\section{Iloczyn skalarny}

\begin{definicja}
Niech $V$ będzie rzeczywistą lub zespoloną przestrzenią liniową. Iloczynem skalarnym nazywamy funkcję
\[
\langle\cdot,\cdot\rangle:V\times V\rightarrow\mathbb K,\qquad \mathbb K=\mathbb R\ \text{lub}\ \mathbb C,
\]
spełniającą następujące warunki:
\begin{enumerate}[label=(I\arabic*)]
    \item liniowość względem pierwszego argumentu:
    \[
    \langle\alpha x+\beta y,z\rangle=\alpha\langle x,z\rangle+\beta\langle y,z\rangle,
    \]
    \item symetrię sprzężoną:
    \[
    \langle x,y\rangle=\overline{\langle y,x\rangle},
    \]
    \item dodatnią określoność:
    \[
    \langle x,x\rangle\geq 0
    \quad\text{oraz}\quad
    \langle x,x\rangle=0\iff x=0.
    \]
\end{enumerate}
\end{definicja}

Dla rzeczywistego iloczynu skalarnego warunek symetrii sprzężonej sprowadza się do
\[
\langle x,y\rangle=\langle y,x\rangle.
\]

\subsection{Norma indukowana przez iloczyn skalarny}
Każdy iloczyn skalarny wyznacza normę:
\[
\|x\|=\sqrt{\langle x,x\rangle}.
\]
Kieruje to nas przez następujący ciąg:
\[
\boxed{\text{iloczyn skalarny}\Longrightarrow\text{norma}\Longrightarrow\text{metryka}.}
\]

\begin{twierdzenie}[Nierówność Cauchy'ego--Schwarza]
Dla wszystkich $x,y\in V$ zachodzi
\[
|\langle x,y\rangle|\leq\|x\|\|y\|.
\]
Równość zachodzi wtedy i tylko wtedy, gdy $x$ i $y$ są liniowo zależne.
\end{twierdzenie}

\subsection{Przykłady iloczynów skalarnych}
\begin{przyklad}
W $\mathbb R^n$:
\[
\langle x,y\rangle=\sum_{i=1}^n x_i y_i.
\]
Wtedy
\[
\|x\|=\sqrt{\sum_{i=1}^n x_i^2}.
\]
\end{przyklad}

\begin{przyklad}
W $\mathbb C^n$, przy przyjętej liniowości względem pierwszego argumentu:
\[
\langle x,y\rangle=\sum_{i=1}^n x_i\overline{y_i}.
\]
\end{przyklad}

\begin{przyklad}
W przestrzeni $\ell^2$:
\[
\langle x,y\rangle=\sum_{n=1}^{\infty}x_n\overline{y_n}.
\]
\end{przyklad}

\begin{przyklad}
W przestrzeni funkcji, na przykład w $L^2[a,b]$:
\[
\langle f,g\rangle=\int_a^b f(t)\overline{g(t)}\,dt.
\]
\end{przyklad}

\subsection{Ortogonalność i twierdzenie Pitagorasa}
Wektory $x,y$ nazywamy ortogonalnymi, jeśli
\[
\langle x,y\rangle=0.
\]
Piszemy wtedy $x\perp y$.

Jeżeli $x\perp y$, to
\[
\|x+y\|^2=\|x\|^2+\|y\|^2.
\]
Jest to twierdzenie Pitagorasa dla normy indukowanej przez iloczyn skalarny.

\subsection{Iloczyn skalarny a kąt między elementami}
W rzeczywistej przestrzeni z iloczynem skalarnym kąt $\varphi\in[0,\pi]$ między niezerowymi elementami $x$ i $y$ określa wzór
\[
\cos\varphi
=\frac{\langle x,y\rangle}{\|x\|\|y\|}.
\]
Z nierówności Cauchy'ego--Schwarza wynika, że iloraz po prawej stronie należy do przedziału $[-1,1]$, więc kąt jest poprawnie określony.

Zatem iloczyn skalarny nie tylko mierzy współliniowość elementów, ale także ich wzajemne ułożenie: w~przestrzeni rzeczywistej
\[
\langle x,y\rangle=\|x\|\|y\|\cos\varphi.
\]
W szczególności $\langle x,y\rangle>0$ oznacza kąt ostry, $\langle x,y\rangle<0$ -- kąt rozwarty, a~$\langle x,y\rangle=0$ -- ortogonalność.

W zespolonej przestrzeni Hilberta iloczyn skalarny może być liczbą zespoloną, dlatego nie można go bezpośrednio potraktować jako iloczynu długości i~kosinusa kąta. Wtedy używa się wzoru
\[
\cos\varphi
=\frac{\operatorname{Re}\langle x,y\rangle}{\|x\|\|y\|},
\]
przy czym $\varphi\in[0,\pi]$ jest kątem wyznaczanym przez część rzeczywistą iloczynu skalarnego.
Wartość $\langle x,y\rangle=0$ nadal oznacza, że elementy są ortogonalne.

\section{Przestrzeń Hilberta}

\begin{definicja}
Przestrzenią Hilberta nazywamy przestrzeń liniową z iloczynem skalarnym, która jest zupełna względem normy indukowanej przez ten iloczyn skalarny.
\end{definicja}

Zatem każda przestrzeń Hilberta jest przestrzenią Banacha, ale nie każda przestrzeń Banacha jest przestrzenią Hilberta. Na przykład normy $\|\cdot\|_1$ i $\|\cdot\|_\infty$ nie są indukowane przez iloczyn skalarny.

\subsection{Przykłady przestrzeni Hilberta}
\begin{itemize}
    \item $\mathbb R^n$ z iloczynem skalarnym
    \[
    \langle x,y\rangle=\sum_{i=1}^n x_i y_i,
    \]
    \item $\mathbb C^n$ z iloczynem skalarnym
    \[
    \langle x,y\rangle=\sum_{i=1}^n x_i\overline{y_i},
    \]
    \item $\ell^2$ z iloczynem skalarnym
    \[
    \langle x,y\rangle=\sum_{n=1}^{\infty}x_n\overline{y_n},
    \]
    \item $L^2[a,b]$ z iloczynem skalarnym
    \[
    \langle f,g\rangle=\int_a^b f(t)\overline{g(t)}\,dt.
    \]
\end{itemize}

\subsection{Związek między typami przestrzeni}
Schemat zależności można zapisać jako:
\[
\boxed{
\begin{gathered}
\text{przestrzeń Hilberta}
\subset \text{przestrzeń Banacha}
\subset \text{przestrzeń unormowana}\\
\subset \text{przestrzeń metryczna}.
\end{gathered}
}
\]
Każda przestrzeń Hilberta jest więc jednocześnie przestrzenią Banacha, unormowaną i metryczną. Implikacje odwrotne nie zachodzą na ogół.

\section{Najważniejsze różnice}
\begin{center}
\begin{tabular}{|p{0.25\linewidth}|p{0.32\linewidth}|p{0.33\linewidth}|}
\hline
\textbf{Typ przestrzeni} & \textbf{Co jest dane?} & \textbf{Dodatkowy warunek} \\
\hline
Metryczna & Zbiór i metryka $d$ & Odległość spełnia trzy aksjomaty metryki \\
\hline
Unormowana & Przestrzeń liniowa i norma $\|\cdot\|$ & Norma spełnia trzy aksjomaty; metryka dana jest przez $d(x,y)=\|x-y\|$ \\
\hline
Banacha & Przestrzeń unormowana & Zupełność względem metryki indukowanej przez normę \\
\hline
Przestrzeń z iloczynem skalarnym & Przestrzeń liniowa i iloczyn skalarny & Norma dana przez $\|x\|=\sqrt{\langle x,x\rangle}$ \\
\hline
Hilberta & Przestrzeń z iloczynem skalarnym & Zupełność względem normy indukowanej przez iloczyn skalarny \\
\hline
\end{tabular}
\end{center}

\end{document}
